%Pour toute question sur l'utilisation de cette classe, écrivez à jerome.soucy(at)ulaval.ca
% Nécessite la classe Evaluation.sty dans le même dossier ou dans le PATH de votre système
% Nécessite le fichier logo_ul.pdf dans le même dossier un dans le PATH de votre système
\def\Ponderation{55}
\def\SigleNumero{STT-1900}
\def\NomCours{Méthodes statistiques pour l'ingénierie}
\def\DateEvaluation{30 avril 2026}
\def\HeureEvaluation{18 h 30 à 21 h 20}
\def\Session{Hiver 2026}
\def\NomEvaluation{Examen 2}

% Ne rien mettre entre les accolades si l'enseignant (ou l'enseignante) est aussi le coordonnateur (ou la coordonnatrice)
\def\Coordonnatrice{}
\def\Coordonnateur{}

% Ne rien mettre entre les accolades s'il y a plus d'une section
\def\Enseignant{Julien Miron}
\def\Enseignante{}

% Ne rien mettre entre les accolades s'il s'agit d'un cours à section unique
% Mettre le nom de chacun des enseignants et enseignantes entre les accolades
% Une case à cocher sera générée pour chacune des sections
\def\SectionA{}
\def\SectionB{}
\def\SectionC{}
\def\SectionD{}

\newcommand{\affichepoints}{}     % ← AVEC points
% \renewcommand{\affichepoints}{on} % ← SANS points (non vide)

\def\Directives{
\begin{enumerate}
\item Matériel autorisé:
\begin{itemize}
\item Calculatrice scientifique {\bf autorisée par la faculté des sciences et de génie};
\item Votre aide-mémoire d'une feuille recto-verso, écrit à la main.
\end{itemize}
\item Assurez-vous que cette évaluation comporte \numquestions ~exercices répartis sur \numpages{}~pages, pour un total de \numpoints{}~points.\
\item \textbf{Veuillez ne pas dégrafer votre examen.}
\item Vous devez\textbf{ justifier toutes vos réponses} par une démarche claire, sauf pour les questions à choix de réponse.\
\item Pour les questions à choix de réponse et les vrais ou faux, vous pouvez mettre un X, un crochet ($\checkmark$) ou remplir la case.\
\item Vous devez laisser tout ce qui ne servira pas durant l'examen à l'avant de la salle AVANT de rejoindre votre place.\
\item Vous n'avez droit à AUCUN appareil électronique en votre possession durant l'examen.
\item Une fois assis à votre place, vous devez demandez l'autorisation pour vous lever, même si l'examen n'est pas encore commencé.
\item Dans les 5 dernières minutes de l'examen, il sera interdit de se lever avant que TOUTES les copies n'aient été ramassées.
\end{enumerate}
\begin{itemize}
\item[$\rightarrow$] Le non-respect des consignes 7-8-9 entraînera une pénalité de 10\%.
\end{itemize}
}

\providecommand{\Gradescope}{}


\documentclass{Evaluation}
\usepackage{multicol,graphicx}

\noprintanswers
\begin{document}


%\pageblanche
%%%%%%%%%%%%%%%% QUESTION %%%%%%%%%%%%%%%%
\begin{questions}

\pageblanche
\question\label{QCMicTests}
Soit $X_1,\hdots,X_{16}$ des réalisations indépendantes d'une loi normale d'espérance $\mu$ et de variance $\sigma^2=16$. L'intervalle de confiance de niveau $95\%$ pour $\mu$ est $[0,04;3,96]$.

\begin{parts}
    \part[2] Quelle est la longueur de la marge d'erreur? \textit{Écrivez votre réponse dans la boîte. N'écrivez rien d'autre que le nombre de votre réponse.}

    \boite{2.5}{1}

    \part[2] Quelle serait la valeur de la borne supérieure de l'intervalle si le niveau était plutôt de $99,15\%$? \textit{Écrivez votre réponse dans la boîte. N'écrivez rien d'autre que le nombre de votre réponse.}

    \boite{2.5}{1}

        \part[2] Quelle est la valeur de $\overline{x}$? \textit{Écrivez votre réponse dans la boîte. N'écrivez rien d'autre que le nombre de votre réponse.}

    \boite{2.5}{1}
\end{parts}
\medskip

Un intervalle de confiance au niveau 95\% pour la variance $\sigma^2$ est $[10,11~;~19,11]$. Pour chaque énoncé qui suit, indiquez si l'énoncé est vrai ou faux.

\medskip

%%%%%%%%%%%%%%%% QUESTION %%%%%%%%%%%%%%%%
\question[3]\label{QCManovaReg}
Montrez que l'ensemble vide $\emptyset$ est à la fois mutuellement exclusif et indépendant de tout événement $A$ pour lequel $0<P(A)<1$.

\vspace{4cm}


\question[3] Soit la table d'analyse de la variance suivante obtenue à l'aide de \texttt{Rcmdr}:
\begin{verbatim}
>summary(ANOVAModel)
            Sum Sq   Df   Mean Sq   F value      Pr(>F)    
Cantons     926.12    4    231.53    238.72    0.000001 
Residuals   446.55   48      0.97
\end{verbatim}

À partir de la table d'analyse de la variance, quelle affirmation parmi les suivantes est \textbf{vraie}?



\begin{checkboxes}
\choice L'estimation de la moyenne du facteur \texttt{Cantons} est égale à $231,53$. 
\choice Nous avons 49 observations dans le jeu de données.
\choice La variance du facteur \texttt{Cantons} est environ 2 fois plus grande que la variance du facteur \texttt{Residuals}.
\correctchoice Au seuil $\alpha=0,05$, nous pouvons affirmer que les moyennes de la variable réponse à chaque niveaux du facteur \texttt{Cantons} ne sont pas toutes égales.
\choice Au seuil $\alpha=0,05$, nous pouvons affirmer que les variances de la variable réponse à chaque niveaux du facteur \texttt{Cantons} ne sont pas toutes égales.
\end{checkboxes}

\pageblanche

%%%%%%%%%%%%%%%% QUESTION %%%%%%%%%%%%%%%%
\question[3] Soit $X_1,\hdots,X_n$ qui proviennent d'une loi normale d'espérance $\mu_X$ et de variance $\sigma^2_X$ et $Y_1,\hdots,Y_n$ qui proviennent d'une loi normale d'espérance $\mu_Y$ et de variance $\sigma^2_Y$. On procède à un test d'égalité des variances avec $H_0:\sigma^2_X=\sigma^2_Y$ et $H_1:\sigma^2_X>\sigma^2_Y$, au seuil $\alpha$, en utilisant la statistique $\displaystyle F_0=\frac{S_Y^2}{S_X^2}$.

\begin{checkboxes}
    \choice Il s'agit d'un test bilatéral.
    \choice Il s'agit d'un test unilatéral à droite.
    \choice Il s'agit d'un test unilatéral à gauche.
    \choice Il nous manque d'information.
\end{checkboxes}

\vspace{0.5cm}
\question[3] 
On définit la fonction de survie $S_X$ d'une variable aléatoire par $S_X(x)=P(X>x)$ et le taux de défaillance $h_X$ par $h_X(x)=\displaystyle\frac{f_X(x)}{S_X(x)}$, où $f_X$ est la fonction de densité de $X$.
Soit la variable aléatoire $X$ dont la fonction de répartition est

\begin{align*}
    F_X(x) = \left\{
			\begin{array}{ll}
			0  & \mbox{ si } x\leq 0, \\
			1-e^{-3x} & \mbox{ si } x> 0.
			\end{array}
			\right. 
\end{align*}
Que vaut $h_X$?

\begin{checkboxes}
    \choice $h_X(x)=3e^{-3x}$ pour $x> 0$ et $0$ sinon.
    \choice $h_X(x)=e^{-3x}$ pour $x> 0$ et $0$ sinon.
    \choice $h_X(x)=3$ pour $x> 0$ et $0$ sinon.
    \correctchoice $h_X(x)=1$ pour $x>0$ et $0$ sinon.
\end{checkboxes}

\vspace{0.5cm}
\question[3] Soit $W\sim\mathcal{N}(1275; 49)$. Laquelle des probabilités suivantes est la plus grande?

\begin{checkboxes}
    \choice $P(W>1282)$
    \choice $P(W<1268)$
    \correctchoice $\displaystyle P\left(\frac{W-1275}{7}>0,5\right)$
    \choice $\displaystyle P\left(\frac{W-1275}{7}>1\right)$
\end{checkboxes}

\vspace{0.5cm}

\question[3] Parmi les quatre façons de générer $X$ et $Y$, laquelle donnera $\rho_{X,Y}=1$?

\begin{checkboxes}
    \choice $X\sim\mathcal{N}(0,1)$ et $Y\sim\mathcal{N}(0,1)$
    \correctchoice $X\sim\mathcal{N}(0,1)$ et $Y=X$
    \choice $X\sim\mathcal{N}(0,1)$ et $Y=X^2$
    \choice $X\sim\mathcal{N}(0,1)$ et $Y=X+\epsilon$ où $\epsilon\sim\mathcal{N}(0,1)$
\end{checkboxes}

\pageblanche

\question Dites si les énoncés suivants sont vrais ou faux.

\begin{parts}
    \part[1] Soit $x_1,\hdots,x_n$. Si la médiane échantillonnale est $25$, nous sommes certains que la valeur de l'observation du milieu est $25$.

    \begin{oneparcheckboxes}
        \choice Vrai
        \correctchoice Faux
    \end{oneparcheckboxes}
    
    \vspace{0.5cm}
        \part[1] Soit $x_1,\hdots,x_n$. Si la moyenne échantillonnale est $25$, environ $50\%$ des observations sont plus grandes que $25$ et environ $50\%$ sont plus petites que $25$.

    \begin{oneparcheckboxes}
        \choice Vrai
        \correctchoice Faux
    \end{oneparcheckboxes}
    
    \vspace{0.5cm}
        \part[1] Soit $x_1,\hdots,x_n$. Si $Q_1=12$, nous savons que le minimum de l'échantillon est 12.

    \begin{oneparcheckboxes}
        \choice Vrai
        \correctchoice Faux
    \end{oneparcheckboxes}
    
    \vspace{0.5cm}

        \part[1] Soit $x_1,\hdots,x_n$. Si la médiane échantillonnale est $25$, nous sommes certains que la valeur de l'observation du milieu est $25$.

    \begin{oneparcheckboxes}
        \choice Vrai
        \correctchoice Faux
    \end{oneparcheckboxes}
    
    \vspace{0.5cm}

            \part[1] Soit $x_1,\hdots,x_n$. Nous devons changer au moins $50\%$ des observations pour que la moyenne échantillonnale change.

    \begin{oneparcheckboxes}
        \choice Vrai
        \correctchoice Faux
    \end{oneparcheckboxes}
    
    \vspace{0.5cm}

            \part[1] Soit $x_1,\hdots,x_n$ et la médiane échantillonnale $\Tilde{x}$. Si on remplace $x_{(n)}$ par $10x_{(n)}$ (c'est-à-dire une valeur 10 fois plus grande), la nouvelle médiane est $10\Tilde{x}$.

    \begin{oneparcheckboxes}
        \choice Vrai
        \correctchoice Faux
    \end{oneparcheckboxes}
    
    \vspace{0.5cm}
\end{parts}

\question[5] Soit $X_1,\hdots,X_n$ qui indépendants et identiquement distribués selon une loi quelconque $F_X$. 
On veut estimer la médiane $\Tilde{\mu}$ de $F_X$ à l'aide de $\overline{X}$.
Calculez le biais de $\overline{X}$ pour $\Tilde{\mu}$.
\pageblanche
 \question Soit $X_1,\hdots, X_{14}$, des variables aléatoires indépendantes et identiquement distribuées selon une loi normale d’espérance $\mu$ et de variance $\sigma^2$. 
Nous désirons construire l'intervalle de confiance approprié de niveau $95\%$ pour $\sigma^2$. Nous notons cet intervalle $[C_1;C_2]$.
\begin{parts}
    \part[3] Que vaut $\mathbb{E}[C_2]$?

    \begin{solution}
        Puisque $\sigma^2 = 400$ est connu et que les $X_i$ suivent une loi normale, l'intervalle de confiance à 95\% est:
        $$\left[\bar{X}-z_{0{,}025}\frac{\sigma}{\sqrt{n}}\;;\;\bar{X}+z_{0{,}025}\frac{\sigma}{\sqrt{n}}\right]=\left[\bar{X}-1{,}96\times\frac{20}{10}\;;\;\bar{X}+1{,}96\times\frac{20}{10}\right]$$
        Donc $C_1=\bar{X}-3{,}92$ et $C_2=\bar{X}+3{,}92$.
        \begin{align*}
            E(C_2)&=E(\bar{X}+3{,}92)\\
            &= E(\bar{X})+3{,}92\\
            &= \mu + 3{,}92
        \end{align*}
    \end{solution}
    \vspace{7cm}
    \part[3] Que vaut $\text{Var}[C_2]$?

    \begin{solution}
        \begin{align*}
            \text{Var}(C_2)&=\text{Var}(\bar{X}+3{,}92)\\
            &= \text{Var}(\bar{X})\\
            &= \frac{\sigma^2}{n}=\frac{400}{100}=4
        \end{align*}
    \end{solution}
    \vspace{7cm}

    \part[4]Quelle est la probabilité que $C_2$ soit supérieure ou égale à $\sigma^2$? 

    \begin{solution}
        \begin{align*}
            P(C_2\leq \mu)&=P(\bar{X}+3{,}92\leq \mu)\\
            &= P\left(\frac{\bar{X}-\mu}{\sigma/\sqrt{n}}\leq \frac{-3{,}92}{20/10}\right)\\
            &= P(Z\leq -1{,}96)\\
            &= 0{,}025
        \end{align*}
    \end{solution}




\end{parts}

\pageblanche

\question[8] Soit $X_1,\hdots, X_{n}$ qui proviennent d'une loi normale d'espérance $\mu$ et de variance $\sigma^2=25$. 
On fait le test d'hypothèses pour lequel $H_0:\mu=5$ et $H_1:\mu>5$ au seuil $\alpha=2,5\%$.
Pour quelles valeurs de $n$ rejette-t-on $H_0$ si $\overline{x}=5,1$?

\pageblanche
\question
Une entreprise fabrique des palettes de chocolat et s'intéresse à la quantité moyenne de sucre dans une palette. Elle choisit au hasard 16 palettes. Dans l'échantillon, la quantité moyenne de sucre par palette est de $5,05 \,g$ et la variance échantillonnale est de $0,011 \, g^2$. Sachant que $X$, la quantité de sucre dans une palette, est une variable aléatoire qui suit une loi normale d'espérance $\mu$ et de variance $\sigma^2=0,01$, elle fait le test suivant, au seuil $\alpha=2,5\%$,
\begin{align*}
	H_0 &: \mu = 5 \cr
	H_1 &: \mu > 5.
\end{align*}

\begin{parts}
\part[2] Quelle est la valeur critique du test en grammes, c'est-à-dire, à partir de quelle valeur de $\overline{X}$ rejette-t-on~$H_0$?

\begin{oneparcheckboxes}
    \choice 1,960
    \choice 2,131
    \choice 2,000
    \CorrectChoice 5,049
    \choice 5,053 
\end{oneparcheckboxes} 

\part[2] Quelle est la valeur de la statistique du test sous $H_0$, arrondie à 3 décimales?

\begin{oneparcheckboxes}
    \choice 1,960
    \choice 2,131
    \CorrectChoice 2,000
    \choice 5,049
    \choice 5,053 
\end{oneparcheckboxes} 

\begin{solution}
Puisque $\sigma^2 = 4$ est connue et $X \sim \text{Normale}$, on utilise le test $Z$. La statistique de test est:
\[
z_0 = \frac{\bar{x} - \mu_0}{\sigma/\sqrt{n}} = \frac{5{,}4 - 5}{2/4} = \frac{0{,}4}{0{,}5} = 0{,}800.
\]
\end{solution}

\part[2] Quel est le seuil observé du test?

	\begin{checkboxes}
	    \choice $\alpha^*=0,025$.
		\CorrectChoice $\alpha^*=0,02275$.
		\choice $\alpha^* \in\ ]0,025;0,05[$ .
		\choice $\alpha^* \in\ ]0,05;0,1[$.
	\end{checkboxes}

\begin{solution}
Le seuil observé est $0{,}212 > 0{,}05 = \alpha$. On ne rejette pas $H_0$. Au seuil $\alpha = 5\%$, la quantité moyenne de sucre dans une palette n'est pas significativement supérieure à $5\,g$.
\end{solution}

\part[3] Si en réalité $\mu=5,1$ et que la règle de décision est de rejeter $H_0$ si $\displaystyle\frac{\overline{X}-5}{0,025}>1,96$, quelle est l'expression de la puissance du test? 

\begin{checkboxes}
    \choice $P\left(\displaystyle \frac{\overline{X}-5,1}{0,025}>1,96\right)$
    \choice $P\left(\displaystyle \frac{\overline{X}-5,1}{0,025}>-2,04\right)$
    \choice $P\left(\displaystyle \frac{\overline{X}-5}{0,025}>1,96\right)$
    \choice $P\left(\displaystyle \frac{\overline{X}-5}{0,025}>-2,04\right)$
\end{checkboxes}

\end{parts}
%%%%%%%%%%%%%%%% FIN DES QUESTIONS %%%%%%%%%%%%%%%%

\pageblanche
\question 
Un ingénieur veut vérifier s'il existe une différence de rigidité d'un certain type de disque, selon le pourcentage de \textit{Pouliovachium} qui le compose.

Il compare donc la rigidité de 8 pourcentages différents (10\%, 20\%, 25\%, 30\%, 35\%, 40\%, 45\% et 50\%) en produisant 4 cylindres pour chaque pourcentage.
La table d'ANOVA suivante est produite par le logiciel utilisé par l'ingénieur. 
Dans celle-ci, certaines quantités ont été remplacées par \texttt{AAA, BBB, CCC} et \texttt{ ???}. Nous considérons que les postulats du modèle sont respectés.
\begin{verbatim}
ANOVA

                     SS      ddl      MS      F0           
pouliovachium    1154.3      ???     ???     AAA
Residuals         141.4      BBB     ???         
Total               CCC      ???       
\end{verbatim}

\begin{comment}
\begin{verbatim}
Analysis of Variance Table

Response: Y
              Df   Sum Sq   Mean Sq   F value        Pr(>F)    
pouliovachium  7   1154.3    164.90     27.99    4.6e-10 ***
Residuals     24    141.4      5.89                    
Total         31   1295.7
\end{verbatim}
\end{comment}

\begin{parts}
    \part[2] Que vaut \texttt{AAA}?
\begin{checkboxes}
    \choice 164,9
    \choice 8,163
    \correctchoice 27,989
    \choice 1154,3
    \choice Il n'est pas possible de trouver la valeur de \texttt{AAA} avec les informations données.
\end{checkboxes}

\begin{solution}
$\texttt{AAA} = F = \dfrac{MS_{\text{trt}}}{MS_E} = \dfrac{278{,}85/5}{343{,}20/24} = \dfrac{55{,}77}{14{,}30} = 3{,}90$, car $\text{Dl}_{\text{trt}} = a-1 = 6-1 = 5$ et $\text{Dl}_E = N-a = 30-6 = 24$.
\end{solution}

\vspace{0.5cm}
    \part[2] Que vaut \texttt{BBB}?
\begin{checkboxes}
    \choice 3
    \choice 4
    \correctchoice 7
    \choice 8
    \choice 31
    \choice 32
    \choice Il n'est pas possible de trouver la valeur de \texttt{BBB} avec les informations données.
\end{checkboxes}

\begin{solution}
$\texttt{BBB} = \text{Dl}_E = N - a = 30 - 6 = 24$.
\end{solution}

\vspace{0.5cm}
    \part[2] Que vaut \texttt{CCC}?
\begin{checkboxes}
    \choice 8,163
    \correctchoice 1295,7
    \choice 1012,9
    \choice 0,122
    \choice Il n'est pas possible de trouver la valeur de \texttt{CCC} avec les informations données.
\end{checkboxes}

\begin{solution}
$\texttt{CCC} = SS_{\text{total}} = SS_{\text{trt}} + SS_E = 278{,}85 + 343{,}20 = 622{,}05$.
\end{solution}

\vspace{0.5cm}
\part[3] Au seuil de 5\%, quelle conclusion pouvons-nous tirer de la table d'ANOVA? Une seule réponse est bonne.
\begin{checkboxes}
    \choice Il n'est pas possible de tirer de conclusion car nous ne connaissons pas les valeurs remplacées par \texttt{???}.
    \choice Il faudrait répéter l'expérience plus de fois pour pouvoir conclure.
    \choice Toutes les moyennes sont significativement différentes.
    \correctchoice La rigidité moyenne des disques d'au moins un pourcentage est significativement différente de celle d'au moins un autre pourcentage.
    \choice Toutes les moyennes sont significativement égales.
\end{checkboxes}

\begin{solution}
La valeur-$p = 0{,}01 < 0{,}05 = \alpha$, donc on rejette $H_0$. La rigidité moyenne des cylindres d'au moins un pourcentage est significativement différente de celle d'au moins un autre pourcentage.
\end{solution}

\end{parts}
\pageblanche

\question[10] Soit le tableau suivant dans lequel certaines informations ont été remplacées par $AAA$ et $BBB$:
\begin{align*}\renewcommand{\arraystretch}{1.4}
   \begin{array}{|c|ccccc|}
   \hline
   i & 1 & 2 & 3& 4& 5\\ \hline
x_i & -3 & -2 & 0 & 1,5 & 3,5   \\  
y_i & -1.1 & AAA & 1,1 & 2 & 2,5 \\ 
x_iy_i & 3,3 & BBB & 0 & 3 & 8,75 \\ 
x_i^2 & 9 & 4 & 0 & 2,25 & 12,25\\ \hline
\end{array} 
\end{align*}

Sachant que $\hat{\beta}_0 = 1$, que vaut $\hat{\beta}_1$? 


\vfill

$\hat{\beta}_1$=\boite{2.5}{1.5}
\vspace{1cm}
\end{questions}

\end{document} 
